\documentclass[11pt]{article}
\usepackage[margin=0.8in]{geometry}
\usepackage[T1]{fontenc}
\usepackage{lmodern,amsmath,amssymb,amsthm,booktabs,array}
\usepackage[colorlinks=true,linkcolor=blue,urlcolor=blue,citecolor=blue]{hyperref}
\hypersetup{pdftitle={Two conjecture proofs and an infinite prime family},pdfauthor={Blake Stone},pdfsubject={Exact source correspondence, complete proofs, and Lean verification}}
\setlength{\parindent}{0pt}
\setlength{\parskip}{0.38em}
\newtheorem{theorem}{Theorem}[section]
\newtheorem{lemma}[theorem]{Lemma}
\newtheorem{corollary}[theorem]{Corollary}
\newcommand{\N}{\mathbb N}
\newcommand{\Z}{\mathbb Z}
\newcommand{\Q}{\mathbb Q}
\newcommand{\R}{\mathbb R}
\begin{document}
\begin{center}
{\LARGE\bfseries Two conjecture proofs\\[3pt]and an infinite prime family}\par
\vspace{0.5em}
{\large Blake Stone}\par
Proof developed with ChatGPT\par
November 27, 2026
\end{center}

\textbf{Abstract.}
We prove the even-degree factorization clause in OEIS A326925:
every normalized degree-$2k$ numerator from the derivatives of
$(1-x)/(1-x-x^2)$ has exactly two distinct irreducible rational factors,
each of degree $k$. Their explicitly normalized constant terms are
$F_{k-1}$ and $L_{k-1}$. Together with the preceding project's
odd-degree theorem, this completes the source's odd/even factorization
pattern. The proof uses two quadratic extensions, norm obstructions,
and descent of conjugate factors. We also prove the A288520/A288523
replacement-length conjecture through an exact coding of its literal
binary scanner, obtaining a stronger affine recurrence. Finally, for
every odd prime $p\equiv2$ or $3\pmod5$, we classify the generalized
A400939 walk with division by $p$: it divides exactly once, at index
$p$ or $p+1$, and then follows an explicit quadratic formula forever.
Dirichlet's theorem certifies that this last family is infinite.
This package leaves the generalized stopping conjecture unresolved
outside the proved family. Three Lean packages contain complete source-definition
bridges and independently checked proof terms. Bounded prior-work
checks found no earlier proof of these exact assertions; no exhaustive
historical-priority claim is made.

\begin{center}
\begin{tabular}{@{}l>{\raggedright\arraybackslash}p{4.65in}@{}}
\toprule
Source & Precisely established contribution\\
\midrule
A326925 & One new even-degree factorization clause, with the companion
Fibonacci/Lucas constants; previous odd-degree work is reused and credited.\\
A288520/A288523 & One shared conjecture on the lengths of literal replacement iterates.\\
A400939 & An infinite prime family of the generalized stopping conjecture,
including the unique division and the exact permanent tail.\\
\bottomrule
\end{tabular}
\end{center}



\section{A326925: the even-degree factors and their irreducibility}
\label{sec:even}

Kimberling's A326925~\cite{A326925} conjectures that every odd-degree
polynomial in its derivative family is irreducible, and that every
degree-$2k$ member has exactly two irreducible factors of degree $k$.
The odd-degree assertion was proved in the preceding project
package~\cite{previous}. Here we prove the even-degree clause and identify
the two constant terms. This is one new factorization assertion with a
companion constant-term result; the earlier odd-degree theorem is not
counted again.

Let
\[
 D=1-x-x^2,\qquad A=2x+1,\qquad B=x^2+x-1,
\]
and define degree-indexed polynomials by
\[
 P_0=-1,\quad P_1=1-x,\quad P_{N+2}=AP_{N+1}-BP_N.
\]
The exact source correspondence is
\begin{equation}\label{eq:source-derivative9}
 \left(\frac{1-x}{1-x-x^2}\right)^{(n)}
 =\frac{n!P_{n+1}(x)}{D^{n+1}}.
\end{equation}
Indeed, the recurrence implies
$DP_N'+NA P_N=N P_{N+1}$ by induction; the quotient rule then proves
\eqref{eq:source-derivative9}. The previous formal package proves both
identities and the representation-independent quotient rule. Thus
degree $2k$ corresponds to derivative order $2k-1$. This indexing
agrees with the source's displayed polynomials and executable program.

Define
\begin{equation}\label{eq:two-factors9}
\begin{array}{lll}
 U_0=1,&U_1=x,&U_{k+2}=AU_{k+1}-BU_k,\\
 V_0=-1,&V_1=2-x,&V_{k+2}=AV_{k+1}-BV_k.
\end{array}
\end{equation}

\begin{theorem}\label{thm:even9}
For every integer $k\ge1$,
\[
 P_{2k}=U_kV_k,
 \qquad \deg U_k=\deg V_k=k.
\]
The two factors are irreducible and coprime over $\Q$. The polynomial
$U_k$ is irreducible even over $\Q(\sqrt5)$.
Their constant terms are
\[
 U_k(0)=F_{k-1},\qquad V_k(0)=L_{k-1},
\]
where $F_0=0,F_1=1$ and $L_0=2,L_1=1$ are the usual Fibonacci and Lucas
initial values. The canonical reduced numerator of the source
derivative of order $2k-1$ is exactly $(2k-1)!\,U_kV_k$.
\end{theorem}

The first examples are
\[
\begin{aligned}
 P_2&=x(2-x),\\
 P_4&=(1+x^2)(1+4x-x^2),\\
 P_6&=(1+3x+x^3)(3+3x+6x^2-x^3).
\end{aligned}
\]
The main source records these products. The companion arrays
A328610/A328611 exchange and change the signs of the two degree-one
factors. Our specified normalization removes that ambiguity; we do not
identify an unspecified factor-list position or claim their separate
row-sum observations.

\subsection{Closed forms}

Put
\[
 u=\frac{1+\sqrt5}{2},\qquad v=1-u,\qquad K=\Q(u).
\]
Thus $u+v=1$, $uv=-1$, and $u\ne v$. Checking initial values and the
common recurrence gives
\begin{align}
 (u-v)P_N&=v^2(x+u)^N-u^2(x+v)^N,\label{eq:p-closed9}\\
 (u-v)U_k&=u(x+v)^k-v(x+u)^k,\label{eq:u-closed9}\\
 V_k&=-v(x+u)^k-u(x+v)^k.\label{eq:v-closed9}
\end{align}
Multiplying the last two equations yields the first at $N=2k$.
Cancellation of $u-v$ proves $P_{2k}=U_kV_k$. Using $vu^2=-u$ rewrites
the factor formulas as
\begin{align}
 (u-v)U_k&=-v\big((x+u)^k+u^2(x+v)^k\big),\label{eq:u-transform9}\\
 V_k&=-v\big((x+u)^k-u^2(x+v)^k\big).\label{eq:v-transform9}
\end{align}
The two leading coefficients are nonzero because $u^2\ne1$ and
$-u^2\ne1$. Consequently both degrees are $k$.

\subsection{Two irreducibility tools}

\begin{lemma}[A sufficient signed binomial criterion]\label{lem:signed9}
Let $E$ be a field, $n>0$, and $a\in E$. If neither $a$ nor $-a$ is a
$p$th power in $E$ for any prime $p\mid n$, then $Y^n-a$ is irreducible
over $E$.
\end{lemma}

\begin{proof}
Express the hypothesis as $(b^p)^2\ne a^2$ for every such $p$ and $b$.
We induct through the prime factors of $n$. The exponent one is
immediate. Write $n=pm$ with $p$ prime. The prime-degree binomial
criterion makes $Y^p-a$ irreducible. For $\alpha^p=a$ its root field
has degree $p$ and
\[
 N_{E(\alpha)/E}(\alpha)=(-1)^p(-a),\qquad
 N_{E(\alpha)/E}(\alpha)^2=a^2.
\]
If $(b^q)^2=\alpha^2$ in this field for a prime $q\mid m$, taking norms
would give
\[
 \big(N_{E(\alpha)/E}(b)^q\big)^2=a^2,
\]
contradicting the base hypothesis. The induction hypothesis therefore
makes $Y^m-\alpha$ irreducible. The binomial composition criterion
then proves irreducibility of $Y^{pm}-a$.
All these steps, including the norm and root-field dimension, are
formalized in \path{SignedKummer.lean}. The prime and composition
ingredients are existing mathlib theorems~\cite{kummer}. We use this
sufficient criterion without claiming a new general classification.
\end{proof}

\begin{lemma}[Descent of exchanged factors]\label{lem:descent9}
Suppose $F\subseteq E$, $H\in F[X]$, and $H=qr$ in $E[X]$, with
$q,r$ coprime irreducibles. If an $F$-automorphism of $E$ sends $q$
to $r$, then $H$ is irreducible over $F$.
\end{lemma}

\begin{proof}
In a factorization $H=ab$ over $F$, primality makes $q$ divide one
factor, say $a$. Applying the automorphism makes $r$ divide that same
factor. Coprimality gives $qr\mid a$, so the other factor $b$ is a unit.
The Lean lemma only needs a coefficient homomorphism fixing $F$ and
sending $q$ to $r$; its two applications below are automorphisms.
\end{proof}

We also reuse the proved fractional-linear lemma: if $a\ne b$, $c\ne1$,
and $Y^n-c$ is irreducible over a field, then
\[
 (X-a)^n-c(X-b)^n
\]
is irreducible. To see this, a root $x$ of any irreducible factor cannot
equal $b$, and $y=(x-a)/(x-b)$ satisfies $y^n=c$. The root field thus
has degree at least $n$. Since the displayed polynomial has degree
$n$, its chosen irreducible factor has full degree. This is the exact
degree argument already formalized in the previous package.

\subsection{Nonpower obstructions in the two quadratic fields}

\begin{lemma}\label{lem:gold9}
In $K=\Q(u)$ neither $u$ nor $-u$ is a $p$th power for any prime $p$.
Also $u^2$ is not a $p$th power for any odd prime $p$.
\end{lemma}

\begin{proof}
The norm of each of $u,-u$ is $-1$. A square root in $K$ would have
rational norm with square $-1$, which is impossible. If $b^p=u$ for an
odd prime, then $b$ is an algebraic integer, is irrational, and has norm
$-1$. In the positive real embedding, $1<b<u$; hence its integral trace
$b-1/b$ lies strictly between $0$ and $1$. This is impossible.
Replacing $b$ by $-b$ deals with $-u$.

If $b^p=u^2$ for an odd prime, $b$ is an irrational algebraic integer
of norm $1$, with $1<b<2$. Its integral trace is $b+1/b$, strictly
between $2$ and $3$, another contradiction. In these arguments
integrality follows because a positive power is integral, and
irrationality follows because the specified power is irrational.
The included unchanged \path{GoldenUnit.lean} establishes the
quadratic degree, integrality, norms, and trace facts.
\end{proof}

By Lemma~\ref{lem:signed9}, $Y^n-u$ and $Y^n+u$ are therefore
irreducible over $K$ for every $n>0$. The odd-exponent binomial
criterion~\cite{kummer} and Lemma~\ref{lem:gold9} also give irreducibility
of $Y^n-u^2$ and $Y^n+u^2$ for odd $n$.

To cover even exponents in the latter plus-sign family, set $J=K(i)$,
where $i^2=-1$. This has degree two over $K$, since $K$ is a real field.

\begin{lemma}\label{lem:imaginary9}
Neither $iu$ nor $-iu$ is a $p$th power in $J$ for any prime $p$.
\end{lemma}

\begin{proof}
For odd $p$, the equation $b^p=\pm iu$ gives, on taking the norm $J/K$,
\[
 N_{J/K}(b)^p=u^2,
\]
contradicting Lemma~\ref{lem:gold9}.

For $p=2$, write $b=a+ci$ with $a,c\in K$. The equation $b^2=\pm iu$
gives $a^2=c^2$ and $2ac=\pm u$. Thus $a=\pm c$ and $u=\pm2c^2$.
The norm $K/\Q$ then gives
\[
 -1=4N_{K/\Q}(c)^2,
\]
which is impossible. Unique coordinates and all the needed norm
identities are checked in the quadratic-extension module.
\end{proof}

\begin{corollary}\label{cor:plus9}
For every $k>0$, $Y^k+u^2$ is irreducible over $K$.
\end{corollary}

\begin{proof}
The odd exponents were already handled. For $k=2m$,
\[
 Y^{2m}+u^2=(Y^m-iu)(Y^m+iu).
\]
Lemma~\ref{lem:imaginary9} and the signed criterion make both factors
irreducible over $J$. They are coprime, since their constants differ.
Complex conjugation exchanges them and fixes $K$.
Lemma~\ref{lem:descent9} proves the conclusion.
\end{proof}

\subsection{Completion of the factorization theorem}

Corollary~\ref{cor:plus9}, the fractional-linear lemma, and
\eqref{eq:u-transform9} show that $U_k$ is irreducible over $K$ and
hence over $\Q$ for every $k>0$.

For odd $k$, irreducibility of $Y^k-u^2$ and
\eqref{eq:v-transform9} show the same for $V_k$. For $k=2m$, set
\[
 q=(x+u)^m-u(x+v)^m,\qquad
 r=-v\big((x+u)^m+u(x+v)^m\big).
\]
Then $V_{2m}=qr$. The irreducibility of $Y^m-u$ and $Y^m+u$,
followed by the fractional-linear lemma, proves that $q,r$ are
irreducible over $K$. Their base polynomials $X+u,X+v$ are coprime,
and their two constants $u,-u$ differ, so $q,r$ are coprime.
The nontrivial automorphism $\sigma$ of $K/\Q$ interchanges $u,v$, and
\[
 \sigma(q)=(x+v)^m-v(x+u)^m=r
\]
because $uv=-1$. Lemma~\ref{lem:descent9} proves that $V_{2m}$ is
irreducible over $\Q$.

The two expressions in \eqref{eq:u-transform9} and
\eqref{eq:v-transform9} are coprime: their constants $u^2,-u^2$ differ,
and $(X+u)^k,(X+v)^k$ are coprime. Nonzero scalar multiplication
preserves this fact, and coprimality descends to $\Q$.
Thus the two irreducible factors are distinct.

It remains to justify reduction of the source rational function.
Since $D=-(X+u)(X+v)$, equation~\eqref{eq:p-closed9} gives, for $N>0$,
\[
 (u-v)P_N(-u)=-u^2(v-u)^N\ne0,\qquad
 (u-v)P_N(-v)=v^2(u-v)^N\ne0.
\]
Hence $P_N$ is coprime to $D$. No polynomial factor cancels in
\eqref{eq:source-derivative9}. For derivative order $2k-1$, the
denominator $D^{2k}$ is monic, so the canonical numerator is precisely
$(2k-1)!\,U_kV_k$.

Finally, at $x=0$ the two recurrences
\eqref{eq:two-factors9} become additive recurrences. Their first two
relevant values are
\[
 U_1(0)=0,\quad U_2(0)=1,\qquad
 V_1(0)=2,\quad V_2(0)=1.
\]
They are therefore $F_{k-1}$ and $L_{k-1}$, respectively. This
completes Theorem~\ref{thm:even9}. \hfill$\square$

\paragraph{Formal endpoint and historical scope.}
The main endpoint is
\begin{quote}\small
\path{A326925Even.odd_derivative_numerator_factorization}.
\end{quote}
It states the exact canonical numerator, irreducibility and coprimality
of its two factors, and both degrees. The factor constants are the
all-index theorems \path{EvenFactors.fib_constant} and
\path{EvenFactors.lucas_constant}. The package includes all thirteen
Lean modules and the unchanged previous dependencies.
The exact official source exports and targeted literal and semantic
prior-work checks are preserved. No earlier proof of the exact
even-degree assertion was located. This bounded search does not
certify historical priority or editorial acceptance.



\section{A288523: a replacement-length conjecture}
\label{sec:a288523}

Start from the binary word $w_0=00$. At each step, scan the current word
from left to right and make the nonoverlapping replacements
\[
  00\longmapsto0101,\qquad 10\longmapsto010,
\]
copying every unmatched letter. This is the literal operation defining
OEIS A288520~\cite{A288520,A288523}. Kimberling conjectured that the lengths of these words
agree with A288523. The same assertion occurs in both entries and is
counted once here.

\begin{theorem}
Writing $L_n=|w_n|$, one has
\[
  (L_0,L_1,L_2,L_3,L_4)=(2,4,5,8,11)
\]
and
\[
  L_N=L_{N-1}+L_{N-2}+L_{N-3}-2L_{N-4}\qquad(N\ge5).
\]
In fact the shorter affine relation
\[
  L_{n+3}+2=L_{n+1}+2L_n\qquad(n\ge1)
\]
holds.
\end{theorem}

\begin{proof}
Introduce tokens $A,B,C$ with decoded binary words
\[
 d(A)=01,\qquad d(B)=001,\qquad d(C)=0001
\]
and extend $d$ by concatenation. Define a token substitution
\[
 h(A)=B,\qquad h(B)=C,\qquad h(C)=BAA.
\]
For a token word $v$, set
\[
 E_0(v)=01\,d(v)\,01,\qquad E_1(v)=001\,d(v)\,01.
\]
Let $R$ denote one literal replacement pass. The exact boundary identity
\begin{equation}\label{eq:a288523-boundary}
 R(1\,d(v)\,01)=01\,d(h(v))\,01
\end{equation}
holds for every finite token word $v$.

For the empty word, it says $R(101)=0101$. To prove the induction step,
suppose that the leading token is $A$, $B$, or $C$. The scan emits,
respectively, $010$, $0100$, or $0100101$ before reaching the residual
input $1\,d(v')\,01$. Applying the induction hypothesis to that residual
input gives the three finite equalities
\[
 010\,01=01\,001,\qquad
 0100\,01=01\,0001,\qquad
 0100101\,01=01\,0010101.
\]
These are precisely the required codes of $h(A),h(B),h(C)$, so
\eqref{eq:a288523-boundary} follows. In particular, this argument proves
the boundary behavior of the nonoverlapping scanner rather than
assuming that rewriting respects arbitrary binary concatenations.

Applying the boundary identity after the two initial prefixes gives
\begin{equation}\label{eq:a288523-states}
 R(E_0(v))=E_1(h(v)),\qquad
 R(E_1(v))=E_0(AA\,h(v)).
\end{equation}
Since $w_1=0101=E_0(\varnothing)$, induction shows that every actual
iterate after the initial word has one of these encoded forms.

Write $\ell_j(v)=|d(h^j(v))|$. Decoded length is additive under
concatenation, and direct evaluation on the three tokens gives
\[
\begin{array}{c|ccc}
 &\ell_0&\ell_1&\ell_3\\\hline
 A&2&3&7\\
 B&3&4&10\\
 C&4&7&15
\end{array}
\]
Consequently, for every token word,
\begin{equation}\label{eq:a288523-cubic}
 \ell_3(v)=\ell_1(v)+2\ell_0(v).
\end{equation}
The lengths after zero, one, two, and three further scans, computed from
\eqref{eq:a288523-states}, are
\[
\begin{array}{c|cccc}
 &0&1&2&3\\\hline
 E_0(v)&\ell_0+4&\ell_1+5&\ell_2+8&\ell_3+11\\
 E_1(v)&\ell_0+5&\ell_1+8&\ell_2+11&\ell_3+16
\end{array}
\]
In both rows, \eqref{eq:a288523-cubic} makes the last entry plus $2$
equal to the second entry plus twice the first. This proves
$L_{n+3}+2=L_{n+1}+2L_n$ for every $n\ge1$. Subtracting neighboring
instances yields
\[
 L_{n+4}+2L_n=L_{n+3}+L_{n+2}+L_{n+1}\qquad(n\ge1),
\]
which is the source recurrence with $N=n+4$. The first five literal
iterates supply the stated initial lengths. Recurrence uniqueness then
gives equality with A288523 at every index.
\end{proof}

The source generating function, an algebraic consequence of the proved
recurrence and initial values, is
\[
 \sum_{n\ge0}L_nx^n=
 \frac{2+2x-x^2-3x^3-2x^4}{1-x-x^2-x^3+2x^4}.
\]
This display is not counted as an additional conjecture resolution or
an additional formal theorem.

\paragraph{Formal correspondence.}
The Lean package defines the scanner directly on lists of Boolean
values, consumes nonoverlapping input matches, and constructs the
literal iterates from $00$. Theorems
\path{Kimberling288523.rewrite_tail} and
\path{Kimberling288523.word_coding} establish the exact connection to
the token coding. The endpoint
\path{Kimberling288523.original_conjecture} identifies every literal
word length with a separately defined integer recurrence, so truncated
natural-number subtraction is not involved. The source recurrence
starts at $N=5$; it is false at $N=4$ and is not asserted there.

\paragraph{Sources and prior work.}
The package preserves exact official OEIS exports for A288523,
revision 11 of April 7, 2020, and A288520, revision 14 of April 6, 2020.
The source review records searches using both engines, the trureturing
index, and the current project index. No earlier proof of this exact
assertion was located; the review does not certify historical priority.
The nearby paper by Dvo\v{r}\'akov\'a, Pelantov\'a, and Shallit,
\emph{On a sequence of Kimberling and its relationship to the
Tribonacci word}, arXiv:2510.11318v5 (2026), concerns A288462 and the
different rule $00\to0101,1\to10$. The trureturing proof of A288476
concerns $00\to0101,1\to011$. Those works illustrate the standard
token-substitution method; no new general method is claimed. The
adjacent position and frequency assertions in A288521 and A288522,
and existence of infinite limiting words, are outside this result's
formal scope.

\section{An infinite prime family for the generalized halving Recam\'an walk}
\label{sec:prime-division-recaman}

OEIS A400939~\cite{A400939}, introduced by Zlatko Damijanic on October 10, 2026,
defines a walk that divides its preceding value by two when that produces
an unvisited integer, and otherwise adds the current index. Its
generalized-divisor conjecture replaces two by an integer that is not a
power of two and predicts eventual permanent addition. We prove an
infinite family of that conjecture, with the stronger conclusion that
exactly one division occurs. The full conjecture and the original
divisor-two occurrence conjecture are not settled here.

For a positive integer $k$, the source rule is
\[
 a_k(1)=1,\qquad
 a_k(n+1)=
 \begin{cases}
  a_k(n)/k,& k\mid a_k(n)\text{ and }
       a_k(n)/k\notin\{a_k(1),\ldots,a_k(n)\},\\
  a_k(n)+n+1,&\text{otherwise},
 \end{cases}
 \qquad n\ge1.
\]
Only the division branch requires an unvisited value. Write
$T_n=n(n+1)/2$.

\paragraph{Theorem.}
Let $p$ be an odd prime with $p\equiv2$ or $3\pmod5$, and put
$q=(p-1)/2$. If $q$ is not a positive triangular number, the unique
division occurs at index $d=p$, and
\[
 a_p(n)=T_n\quad(1\le n<p),\qquad
 a_p(n)=\frac{n(n+1)-p^2-1}{2}\quad(n\ge p).
\]
If $q$ is a positive triangular number, the unique division occurs at
$d=p+1$, and
\[
 a_p(n)=T_n\quad(1\le n\le p),\qquad
 a_p(n)=\frac{n(n+1)-(p+1)^2}{2}\quad(n\ge p+1).
\]
In each case, $p\nmid a_p(n)$ for every $n\ge d$. There are arbitrarily
large primes to which this theorem applies.

\paragraph{Proof.}
If $1\le n$ and $n+1<p$, then $p\nmid T_n$: a contrary divisibility
would imply $p\mid n(n+1)$ although neither positive factor is as large
as $p$. Induction in the literal walk therefore gives
$a_p(n)=T_n$ for $1\le n<p$.

At the next step, the preceding value is
$T_{p-1}=p(p-1)/2=pq$. If $q$ is unvisited, division gives
$a_p(p)=q$, and hence
\[
 2a_p(p)+p^2+1=p(p+1).
\]
If $q$ has appeared, the rule instead gives $a_p(p)=T_p$. The next
proposed quotient is $q+1$. Two positive triangular numbers cannot be
consecutive: for $i>j\ge1$,
\[
 T_i-T_j\ge T_{j+1}-T_j=j+1\ge2.
\]
Thus $q+1$ has not appeared in the triangular prefix, and
$a_p(p+1)=(p+1)/2$. In this case,
\[
 2a_p(p+1)+(p+1)^2=(p+1)(p+2).
\]
This initial-division argument holds for every odd prime.

In the first case let $D=p^2+1$, and in the second let $D=(p+1)^2$.
At the first division index $d$ we have
\[
 2a_p(d)+D=d(d+1),\qquad D\equiv1\pmod p.
\]
Suppose that $2a_p(n)+D=n(n+1)$ for some $n\ge d$. If
$p\mid a_p(n)$, reduction modulo $p$ would give $n(n+1)\equiv1$,
and consequently
\[
 (2n+1)^2=4n(n+1)+1\equiv5\pmod p.
\]
Quadratic reciprocity, since $5\equiv1\pmod4$, identifies the quadratic
character of $5$ modulo $p$ with that of $p$ modulo $5$. The only square
residues modulo $5$ are $0,1,4$, so the residue classes in the theorem
make this impossible. Hence the walk must add $n+1$. The invariant is
preserved because
\[
 2a_p(n+1)+D=n(n+1)+2(n+1)=(n+1)(n+2).
\]
Induction proves the stated tail and its nondivisibility. Together with
the initial dynamics, this proves uniqueness of the division.

Finally, Dirichlet's theorem gives arbitrarily large primes
$p\equiv3\pmod5$. These primes form an infinite subfamily of the
classification. This last step and quadratic reciprocity are credited
classical dependencies.\hfill$\square$

\paragraph{Formalization and source scope.}
The accompanying Lean module \path{A400939/Proof.lean} defines the
literal visited-value walk, explicitly excluding its auxiliary zero
index from the visited set. The declaration
\path{walk_followsRule} proves the exact source-rule bridge;
\path{first_prime_division} checks the first division for all odd
primes; and \path{division_at_prime_iff} checks the triangular test
selecting the index. The main theorem
\path{walk_prime_residue_classes} proves the full prefix, unique
division, permanent exact tail, and tail nondivisibility.
\path{arbitrarily_large_single_division} proves infinitude using
mathlib's existing Dirichlet theorem. Both end theorems depend only on
the standard Lean axioms \path{propext}, \path{Classical.choice},
and \path{Quot.sound}.

The primary statement is pinned to OEIS A400939 revision 58,
November 5, 2026, in OEIS data-export commit
\path{df38d77b83a78d8ab5aaa501dec6e3a3db977f01}.
Bounded exact-identifier and semantic searches of both available web
engines, the current trureturing catalog, and its GitHub issues found
no prior proof of this particular prime-family classification. This
documents a bounded check rather than establishing worldwide priority.
The package is an infinite-family partial result, not a claim that
the entire A400939 generalized conjecture has been resolved.


\section{Verification, source fidelity, and reproducibility}

Each mathematical endpoint is proved from the corresponding literal
source definition or an explicitly proved source bridge. Polynomial
irreducibility refers to the actual canonical reduced numerator, not
merely an uncancelled expression. The replacement theorem starts from
the specified nonoverlapping scanner. The walk theorem retains the
visited-value test on divisions and permits repeats after addition.

The three independent projects use Lean 4.19.0 and mathlib revision
\[
\texttt{c44e0c8ee63ca166450922a373c7409c5d26b00b}.
\]
From the relevant package directory, the reproducible commands are
\begin{verbatim}
lake exe cache get
lake build
\end{verbatim}
The projects are \path{A326925Even}, \path{A288523}, and
\path{A400939}. The first includes thirteen local Lean modules; its
Lake configuration lists all module roots. The other two have a
single proof module. The prior five-module A326925 development is
included unchanged in the new algebra project, with its old
\path{Proof.lean} renamed to \path{OddProof.lean}; it is a credited
dependency rather than a new result.

Author and independent builds used fresh project directories without
copied local proof objects. The repository includes their logs, exact
source hashes, source-definition reviews, and pinned primary exports.
The binary-word endpoints use only \path{propext} and
\path{Quot.sound}; the algebra and walk endpoints use those plus
\path{Classical.choice}. No admission, added axiom, unsafe declaration,
or native-evaluation proof shortcut is used.

The formal dependencies include existing proofs of the prime/odd
binomial criteria, quadratic reciprocity, and Dirichlet's theorem.
The arbitrary-exponent signed helper is a sufficient criterion
proved here for the application, without a claim to new general
binomial theory. Its norm induction is stronger than the imported
odd-exponent argument and handles the required even exponents.

The current OEIS exports, literal-ID searches, semantic searches,
and comparisons with nearby public solutions are documented in each
package's source review. Several unrelated candidates were discarded
after prior proofs were found. The result count here concerns two
recorded conjectural assertions, the stated constant-term companion,
and one infinite-family partial result. The two entries sharing the
word assertion are counted once; generic supporting lemmas and the
previous odd-degree theorem are not counted as new problems solved.
No claim is made of OEIS editorial acceptance, peer review by human
mathematicians, or exhaustive historical novelty.

\begin{thebibliography}{99}
\raggedright
\bibitem{A326925}
C. Kimberling, \emph{OEIS A326925}, October 22, 2019.
\url{https://oeis.org/A326925}.
Official export revision 13, October 31, 2019.

\bibitem{companions}
C. Kimberling, \emph{OEIS A328610 and A328611}.
\url{https://oeis.org/A328610};
\url{https://oeis.org/A328611}.
Official export revisions 12 and 11, October 31, 2019.

\bibitem{previous}
B. Stone, \emph{A326925: irreducibility of every odd-degree derivative
numerator}, previous project package, November 27, 2026.
\url{https://github.com/blakestone-x/math/tree/main/A326925}.
Proof developed with ChatGPT; the package includes its full Lean
development and independent audit.

\bibitem{kummer}
The mathlib contributors, \emph{KummerPolynomial.lean} and
\emph{KummerExtension.lean}, pinned mathlib revision above.
\url{https://github.com/leanprover-community/mathlib4/tree/c44e0c8ee63ca166450922a373c7409c5d26b00b/Mathlib/FieldTheory}.
The odd-exponent theorem
\path{Polynomial.X_pow_sub_C_irreducible_of_odd}
is in \path{KummerExtension.lean}; this corrects the earlier
package's filename citation without altering its mathematics.

\bibitem{A288520}
C. Kimberling, \emph{OEIS A288520}, June 13, 2017.
\url{https://oeis.org/A288520}.
Official export revision 14, April 6, 2020.

\bibitem{A288523}
C. Kimberling, \emph{OEIS A288523}, June 13, 2017.
\url{https://oeis.org/A288523}.
Official export revision 11, April 7, 2020.

\bibitem{nearbyword}
L. Dvo\v{r}\'akov\'a, E. Pelantov\'a, and J. Shallit,
\emph{On a sequence of Kimberling and its relationship to the
Tribonacci word}, arXiv:2510.11318v5, 2026.
\url{https://arxiv.org/abs/2510.11318}.
This concerns A288462 and a different replacement rule.

\bibitem{A400939}
Z. Damijanic, \emph{OEIS A400939}, October 10, 2026.
\url{https://oeis.org/A400939}.
Official export revision 58, November 5, 2026.

\bibitem{quadratic}
The mathlib contributors, \emph{QuadraticReciprocity.lean},
pinned revision above, theorem
\path{ZMod.exists_sq_eq_prime_iff_of_mod_four_eq_one}.
\url{https://github.com/leanprover-community/mathlib4/tree/c44e0c8ee63ca166450922a373c7409c5d26b00b/Mathlib/NumberTheory/LegendreSymbol}.

\bibitem{dirichlet}
The mathlib contributors, \emph{PrimesInAP.lean},
pinned revision above, theorem
\path{Nat.forall_exists_prime_gt_and_eq_mod}.
\url{https://github.com/leanprover-community/mathlib4/tree/c44e0c8ee63ca166450922a373c7409c5d26b00b/Mathlib/NumberTheory/LSeries}.
\end{thebibliography}
\end{document}

